\section{Targeted Languages for Embedded Processor}
There are vast range of languages suitable to program embedded systems, like-- assembly, C, C++, and Java. However, there are differences in many aspects--- ``for example in the abstraction level at which the program is written and with that the granularity of control over the processor''~\cite{Gajski.2009}. Among them, assembly language represent the lowest level of abstraction in software development. That is why it is mentioned by~\cite{Gajski.2009} that the assembly is ``basically a symbolic representation of the processor’s machine code and only minimally abstracts the processors complexity''. It gives control over the processor's internals--- such as controlling the registers and instruction-level operations. That is why, it is also valuable for writing performance-critical or low-level routines. However, assembly is very much dependent on the processor's Instruction Set Architecture (ISA), meaning that each processor family (e.g., ARM, x86, RISC-V) differs in its own unique assembly language.